\section{Procedure}\label{sec:procedure}

%Presents the apparatus and procedure used and includes:
%1) a list of the apparatus;
%2) an apparatus diagram; 3) the technical details of the investigation in the procedure used to acquire data.
%Written in the past tense with specific terminology, this should allow others to replicate the experiment.

The procedure below followed the outline in the PHYS 326 course lab book ~\cite{phys326}.

Two methods were used to find the dipole moment of each magnet: Method 1 used a Gauss meter to measure the magnetic
field strength
and plotting the far-field behaviour, and method 2 used the EMF curve from dropping the magnets through successive
coils.

A ruler was used to measure the axial distance $z$ whilst placing the Gauss meter at regular intervals away from the
magnet.
$B_z$ and their corresponding distance values were recorded in a table and incremented by 5mm until the magnetic
field strength became too weak to measure.
This step was done for both the short and long magnets.

\begin{figure}[htb]
        \centering
        \includegraphics[width=\linewidth,height=0.35\textheight,keepaspectratio]{figures/measure}
        \caption{Method 1: The magnet was placed along the ruler. $B_z$ and corresponding distance measurements were taken.}
        \label{fig:meas}
    \end{figure}

The slope of the line where it was approximately linear was used to derive the magnetic dipole moment using Eq. ~\eqref{eq:Bz}.

As for method 2, a magnet was carefully attached to the top of the string and left to fall (with North polarity
facing down) from rest.
With the coils hooked up to dual input amplifiers, this resulted in EMF which was sampled by the LabPro sensor
hardware.
This data was recorded in LabPro at the highest sampling rate available, and EMF curves for each coil were plotted
against time.
The EMF curves were de-amplified by the gain rating used to reflect the actual response.
This was repeated for both the short and long magnets.

\clearpage


%The experimental apparatus and procedure followed the standard method outlined in the University of Victoria PHYS 326 laboratory manual~\cite{phys326}.

%Table~\ref{tab:params} provides a list of measurements taken.
%Uncertainty was calculated using...


%\begin{table}[htbp]
%\caption{Experimental Parameters}
%\begin{tabular}{lrrr}
%\toprule
%{Quantity} & {Reading} \\
%\midrule
%$R_L$, Inductor resistance & \qty{109.4(1.3)}{\ohm} \\
%L, Inductance & \qty{1.044(0.072)}{\henry} \\
%C, Capacitance & \qty{10.07(0.21)}{\micro\farad}\\
%$R_{FG}$, function generator resistance & \qty{50.00(0.25)}{\ohm} \\
%$R_{res}$, normal & \qty{1133}{\ohm} \\
%$R_{res}$, critical & \qty{485}{\ohm} \\
%$R_{res}$, oscillatory & \qty{20}{\ohm} \\
%\bottomrule
%\end{tabular}\label{tab:params}
%\end{table}